\chapter{Plugin-Verwaltung}\label{kap:plugin-verwaltung}\index{Plugin-Verwaltung}

Dieser Anhang beschreibt, woher \produkt{} seine Plugins nimmt, wie Sie Plugin-Pakete
installieren und entfernen, wie die Signaturprüfung arbeitet, was ein Plugin in seiner
\datei{plugin.json} an Python-Paketen verlangt und wie die Update-Prüfung funktioniert. Die
Bedienung der Plugins selbst beschreibt Kapitel~\ref{kap:plugins}.

% ---------------------------------------------------------------------------
\section{Plugin-Ordner}\label{sec:verwaltung-ordner}\index{Plugin-Ordner}

Ein Plugin ist ein Ordner, der eine Datei \datei{plugin.json} enthält. Die Anwendung durchsucht
mehrere \emph{Plugin-Ordner}; jeder Unterordner eines Plugin-Ordners mit einer
\datei{plugin.json} ist ein Plugin. Die Ordner werden in dieser Reihenfolge durchsucht:

\begin{enumerate}
  \item der \textbf{Ordner für installierte Plugins} (Abschnitt~\ref{sec:verwaltung-installiert}),
  \item die \textbf{weiteren Plugin-Ordner} aus \menu{Plugins,Settings…}
        (\gui{Further plugin folders}),
  \item die Ordner aus der Umgebungsvariablen \code{NFCAD\_PLUGIN\_PATH} (mehrere Ordner durch
        \code{:} getrennt, unter Windows durch \code{;}),
  \item in einer Entwicklungsversion, wenn sonst kein Ordner eingestellt ist, die Ordner
        \datei{python/examples} und \datei{python/plugins} des Quellcodes,
  \item zuletzt der \textbf{Ordner der mitgelieferten Plugins}
        (Abschnitt~\ref{sec:verwaltung-mitgeliefert}).
\end{enumerate}

Jede Plugin-Kennung (\code{id}) darf nur einmal geladen werden. Findet die Anwendung eine
Kennung ein zweites Mal, gilt der zuerst gefundene Ordner:

\begin{itemize}
  \item Ein installiertes Plugin hat Vorrang vor einem gleichnamigen in einem anderen Ordner.
  \item Ein Plugin in einem eigenen Ordner \emph{ersetzt} das mitgelieferte Plugin gleicher
        Kennung – ohne Meldung. So können Sie eine neuere Version eines mitgelieferten Plugins
        verwenden.
  \item Kommt eine Kennung in zwei anderen Ordnern vor, wird nur die erste geladen, und
        \gui{Problems:} im Einstellungsdialog meldet \code{plugin id '…' is used twice}.
\end{itemize}

\menu{Plugins,Reload plugins} sucht alle Ordner neu ab, beendet alle Plugin-Prozesse und startet
die gefundenen Plugins neu.

\subsection{Der Ordner für installierte Plugins}\label{sec:verwaltung-installiert}

\gui{Install plugin…} installiert in den Unterordner \datei{plugins} des Datenordners der
Anwendung, je Plugin in einen Ordner mit seiner Kennung (\datei{plugins/<id>}). Der
Einstellungsdialog zeigt den genauen Pfad unter \gui{Installed plugins (Plugins $\to$ Install
plugin…):}. Typisch ist:

\begin{referenz}{@{}lX@{}}
\toprule
\textbf{System} & \textbf{Ordner} \\
\midrule
\endhead
macOS & \datei{\textasciitilde/Library/Application Support/PhotonPower/Raytatouille Fusion/plugins} \\
Windows & \datei{\%APPDATA\%\textbackslash PhotonPower\textbackslash Raytatouille Fusion\textbackslash plugins} \\
Linux & \datei{\textasciitilde/.local/share/PhotonPower/Raytatouille Fusion/plugins} \\
\bottomrule
\end{referenz}

Die Umgebungsvariable \code{NFCAD\_PLUGIN\_INSTALL\_DIR} legt einen anderen Ordner fest.

\subsection{Mitgelieferte Plugins}\label{sec:verwaltung-mitgeliefert}\index{mitgelieferte Plugins}

Die Plugins \plugin{Raytracer}, \plugin{Workflow}, \plugin{FourierGrating}, \plugin{hp-FEM},
\plugin{Path Tracer} und \plugin{Automotive Lighting} werden mit der Anwendung installiert (Windows-Installer,
macOS-Disk-Image, RPM). Sie liegen neben dem Python-Paket \code{nfcad} der Anwendung:

\begin{referenz}{@{}lX@{}}
\toprule
\textbf{System} & \textbf{Ordner (relativ zum Programm)} \\
\midrule
\endhead
macOS & \datei{RaytatouilleFusion.app/Contents/Resources/python/plugins} \\
Windows & \datei{<Installationsordner>\textbackslash python\textbackslash plugins} \\
Linux & \datei{<prefix>/share/raytatouille-fusion/python/plugins} \\
\bottomrule
\end{referenz}

Mitgelieferte Plugins werden \emph{ohne Signatur} geladen: Sie kommen aus demselben Installer und
liegen im selben schreibgeschützten Ordner wie das Programm; eine Signatur würde nichts bezeugen,
was die Installation nicht schon bezeugt. Als mitgeliefert gelten nur diese festen Orte relativ
zum Programm – weder eine Einstellung noch eine Umgebungsvariable kann einen anderen Ordner dazu
machen. Eine Entwicklungsversion, die direkt aus dem Quellcode läuft, hat keinen solchen Ordner.

Mitgelieferte Plugins lassen sich nicht mit \gui{Uninstall plugin} entfernen; nehmen Sie ein
Plugin, das Sie nicht brauchen, stattdessen mit \menu{Plugins,Show in the side panels} aus den
Seitenleisten (Abschnitt~\ref{sec:plugin-seitenleisten}).

% ---------------------------------------------------------------------------
\section{Plugin-Pakete installieren und entfernen}\label{sec:verwaltung-installieren}\index{nfplugin@\texttt{.nfplugin}}

Plugins werden als \emph{Plugin-Paket} mit der Endung \datei{.nfplugin} weitergegeben – ein
ZIP-Archiv des Plugin-Ordners mit seiner Signaturdatei.

\subsection{Installieren}

\begin{schritte}
  \item Wählen Sie \menu{Plugins,Install plugin…} und im Dateidialog
        (\gui{Plugin packages (*.nfplugin)}) das Paket. Alternativ übergeben Sie die
        \datei{.nfplugin}-Datei beim Start als Datei-Argument (Anhang~\ref{kap:programme}),
        z.\,B.\ indem Sie sie im Dateimanager mit der Anwendung öffnen.
  \item Die Anwendung beendet alle Plugin-Prozesse, prüft das Paket und seine Signatur, entpackt
        es in den Ordner für installierte Plugins und lädt alle Plugins neu.
  \item Ein Meldungsfenster \gui{Install plugin} bestätigt: \code{… installed.} bzw.\ bei einer
        schon installierten Version \code{… updated from 0.1.0 to 0.2.0.}
  \item Fehlen in der Python-Umgebung Pakete, die das Plugin verlangt, nennt die Meldung sie mit
        dem fertigen Befehl: \code{Missing Python packages – install them with: …}.
\end{schritte}

Während ein Plugin-Job läuft, ist keine Installation möglich (\gui{A plugin job is still
running}).

\paragraph{Prüfungen beim Installieren.} Das Paket wird abgelehnt, wenn

\begin{itemize}
  \item es kein lesbares ZIP-Archiv ist (\code{… is not a plugin package}),
  \item es mehr als 5000 Dateien enthält, eine Datei größer als 64~MB ist oder alles zusammen
        entpackt mehr als 256~MB groß wäre (\code{too many files in the package},
        \code{package too large}),
  \item ein Pfad im Archiv absolut ist, \code{..} enthält, einen Laufwerksbuchstaben oder einen
        Rückwärtsschrägstrich hat oder ein symbolischer Link ist
        (\code{invalid path in the package: …}),
  \item die \datei{plugin.json} weder ganz oben noch im einzigen Ordner der obersten Ebene liegt
        (\code{the package has no plugin.json}),
  \item die \datei{plugin.json} fehlerhaft ist (Abschnitt~\ref{sec:verwaltung-manifest}),
  \item die Signatur fehlt oder ungültig ist und der Entwicklermodus aus ist
        (\code{… – only signed plugins can be installed (unsigned ones in developer mode)}).
\end{itemize}

Im Entwicklermodus wird ein unsigniertes Paket installiert; die Meldung ergänzt
\code{Not signed (…) – it only loads in developer mode.}

Eine installierte ältere (oder gleiche) Version wird vollständig ersetzt.

\subsection{Entfernen}

\menu{Plugins,Uninstall plugin} listet alle geladenen Plugins mit Name und Version. Wählen Sie ein
installiertes Plugin und bestätigen Sie die Frage \gui{Remove the plugin …?}. Der Ordner des
Plugins wird gelöscht, die Plugins werden neu geladen; die Statuszeile meldet
\code{… removed}.

Plugins aus anderen Ordnern (eigene Plugin-Ordner, mitgelieferte Plugins) stehen ausgegraut mit
dem Zusatz \gui{– from a plugin folder} in der Liste. Der Tooltip nennt den Ordner, aus dem sie
kommen. Solche Plugins entfernen Sie, indem Sie den Ordner unter \menu{Plugins,Settings…}
austragen oder löschen.

% ---------------------------------------------------------------------------
\section{Die Beschreibungsdatei plugin.json}\label{sec:verwaltung-manifest}\index{plugin.json@\texttt{plugin.json}}

Jedes Plugin beschreibt sich in einer JSON-Datei \datei{plugin.json}, z.\,B.\ der Raytracer:

\begin{lstlisting}
{
  "id": "raytracer",
  "name": "Raytracer",
  "version": "0.1.0",
  "api": 1,
  "entry": "raytracer_plugin:RaytracerPlugin",
  "requires": ["numpy", "raytatouille>=0.9", "rtt_fusion>=0.3.2 (rtt-fusion)",
               "fieldtracing>=0.2 (photonpower-fieldtracing)"],
  "description": "Raytatouille: load optical systems, trace rays, ..."
}
\end{lstlisting}

\begin{referenz}{@{}>{\raggedright\arraybackslash}p{0.2\linewidth}X@{}}
\toprule
\textbf{Feld} & \textbf{Bedeutung} \\
\midrule
\endhead
\code{id} & Kennung, Pflicht. Nur Kleinbuchstaben, Ziffern, \code{\_}, \code{-} und \code{.},
  beginnend mit Buchstabe oder Ziffer, höchstens 64 Zeichen. Benennt den Ordner eines
  installierten Plugins und die Parameter des Plugins im Dokument. \\
\code{name} & Anzeigename (ohne Angabe: die Kennung). \\
\code{version} & Version, z.\,B.\ \code{0.1.0}; für Installation und Update-Prüfung. \\
\code{api} & Version der Plugin-Schnittstelle; diese Ausgabe unterstützt nur \code{1}. Ein
  anderer Wert führt zu \code{plugin API … is not supported (this version: 1)}. \\
\code{entry} & Pflicht: Python-Modul und Klasse des Plugins, \code{modul:Klasse}. \\
\code{requires} & Benötigte Python-Pakete (Abschnitt~\ref{sec:verwaltung-requires}). \\
\code{description} & Beschreibung; erscheint im Plugin-Dock und als Tooltip im Menü. \\
\code{update\_url} & Optional: Adresse für die Update-Prüfung
  (Abschnitt~\ref{sec:verwaltung-updates}). \\
\code{previews} & Optional: Blocktypen, deren Übersichtsbild das Plugin zeichnet, obwohl es sie
  nicht selbst anbietet, z.\,B.\ \code{["optics.*", "fusion.particle"]} beim Workflow-Plugin;
  \code{*} am Ende passt auf jeden Typ mit diesem Anfang
  (Abschnitt~\ref{sec:verwaltung-uebersicht}). \\
\bottomrule
\end{referenz}

Fehlt \code{id} oder \code{entry}, meldet der Einstellungsdialog unter \gui{Problems:}
\code{plugin.json needs 'id' and 'entry'}; eine ungültige Kennung
\code{plugin id '…':} \code{use lower-case letters, digits, '\_', '-', '.'}.

\subsection{Python-Anforderungen}\label{sec:verwaltung-requires}\index{Python-Anforderungen}

Jeder Eintrag in \code{requires} hat die Form

\begin{lstlisting}
importname [>= version | == version] [(pip-name)]
\end{lstlisting}

\begin{itemize}
  \item \textbf{Importname}: der Name, unter dem Python das Paket importiert,
        z.\,B.\ \code{fieldtracing}.
  \item \textbf{Version} (optional): \code{>=} verlangt mindestens, \code{==} genau diese
        Version. Verglichen wird numerisch Teil für Teil; ein Zusatz wie \code{-rc1} zählt
        weniger als die Version ohne Zusatz.
  \item \textbf{pip-Name} in Klammern (optional): der Name des Pakets für \code{pip}, wenn er
        vom Importnamen abweicht, z.\,B.\ \code{fieldtracing>=0.2 (photonpower-fieldtracing)}
        oder \code{rtt\_fusion>=0.2 (rtt-fusion)}.
\end{itemize}

Diese Angaben nutzen \gui{Check environment} (es importiert das Paket unter dem Importnamen und
liest seine Version), \gui{Create environment…} und die Meldung nach einer Installation (beide
installieren mit dem pip-Namen). Pakete mit Versionsangabe setzt der angezeigte \code{pip}-Befehl
in Anführungszeichen, z.\,B.\ \code{"photonpower-fieldtracing>=0.2"}. Ein Eintrag, der nicht in
diese Form passt, erscheint als \code{invalid requirement}. Zusätzlich verlangt die Anwendung
immer \code{numpy}, weil ihr Python-Paket \code{nfcad} es für die Datenübergabe braucht.

Die Python-Version muss 3.10 oder neuer sein.

% ---------------------------------------------------------------------------
\section{Signaturen}\label{sec:verwaltung-signaturen}\index{Signatur}

Plugins sind Python-Code, der mit Ihren Rechten läuft. Deshalb lädt \produkt{} außer den
mitgelieferten nur Plugins, die vom Herausgeber signiert sind. Es gibt genau \emph{einen}
vertrauenswürdigen Herausgeber – den Inhaber des Projekts; sein öffentlicher Schlüssel ist in die
Anwendung eingebaut (im Einstellungsdialog als \gui{Trusted publisher key: ed25519:…}). Es gibt
keine Möglichkeit, zur Laufzeit weitere Schlüssel hinzuzufügen, auch nicht über Einstellungen oder
Umgebungsvariablen; ein neuer Schlüssel braucht eine neue Version der Anwendung. Ist kein
Schlüssel eingebaut, steht dort \gui{Trusted publisher: none built in – signed plugins cannot be
loaded …}.

\subsection{Was wird signiert?}

\begin{itemize}
  \item Verfahren: Ed25519 (RFC 8032) mit der Bibliothek Monocypher; Dateiprüfsummen SHA-256.
  \item Signiert werden \emph{alle} Dateien des Plugin-Ordners einschließlich Unterordnern –
        außer der Datei \datei{SIGNATURE} selbst, \datei{.pyc}-Dateien, Ordnern
        \datei{\_\_pycache\_\_} und Dateien oder Ordnern, deren Name mit einem Punkt beginnt.
  \item Signierter Text: die Zeile \code{nfcad-plugin-signature v1}, gefolgt von einer Zeile
        \textit{SHA-256}~\textit{Pfad} je Datei, nach Pfad sortiert (Pfade relativ mit
        \code{/}).
  \item Die Datei \datei{SIGNATURE} (JSON) neben der \datei{plugin.json} enthält die Dateiliste mit
        Prüfsummen, den öffentlichen Schlüssel des Unterzeichners und die Signatur. Die
        \datei{plugin.json} muss zu den signierten Dateien gehören.
\end{itemize}

\subsection{Was passiert beim Laden?}

Die Anwendung kopiert \emph{genau die signierten Dateien} in einen privaten temporären Ordner,
berechnet dabei die Prüfsummen der geschriebenen Bytes und vergleicht sie mit der Signatur. Das
Plugin läuft dann aus dieser Kopie. Dateien, die nach der Prüfung im Plugin-Ordner hinzukommen oder
geändert werden, werden also nie ausgeführt. Die Kopien werden beim nächsten Neuladen entfernt.

Mögliche Ergebnisse der Prüfung (Anzeige unter \gui{Problems:} im Einstellungsdialog bzw.\ beim
Installieren):

\begin{referenz}{@{}>{\raggedright\arraybackslash}p{0.38\linewidth}X@{}}
\toprule
\textbf{Meldung} & \textbf{Bedeutung} \\
\midrule
\endhead
\code{not signed} & Keine Datei \datei{SIGNATURE}. \\
\code{signed with an unknown key (…)} & Signiert, aber nicht vom eingebauten Herausgeber. \\
\code{modified after signing: …} & Eine signierte Datei wurde geändert. \\
\code{signed file missing: …} & Eine signierte Datei fehlt. \\
\code{plugin.json is not signed} & Die Beschreibungsdatei gehört nicht zu den signierten
  Dateien. \\
\code{SIGNATURE is malformed}, \code{cannot read SIGNATURE} & Signaturdatei beschädigt bzw.\
  größer als 4~MB. \\
\code{invalid path in SIGNATURE: …} & Ein Pfad verlässt den Plugin-Ordner. \\
\code{signature does not match} & Die Signatur passt nicht zu den Dateien. \\
\bottomrule
\end{referenz}

Ein solches Plugin wird nicht geladen:
\code{…: not signed – not loaded (only signed plugins; unsigned ones in developer mode)}.

\begin{warnung}
Eine Signatur beweist, \emph{wer} ein Plugin geschrieben hat und dass es unverändert ist. Sie
beweist nicht, dass es harmlos ist – Plugins laufen ohne Sandbox mit Ihren Rechten.
\end{warnung}

\subsection{Der Entwicklermodus}\label{sec:verwaltung-entwicklermodus}\index{Entwicklermodus}

Wer ein eigenes Plugin schreibt, kann es nicht mit dem Schlüssel des Herausgebers signieren. Dafür
gibt es den Entwicklermodus:

\begin{itemize}
  \item Einschalten: \menu{Plugins,Settings…} $\to$ \gui{Developer mode: also load unsigned
        plugins} oder Umgebungsvariable \code{NFCAD\_DEVELOPER\_MODE=1} (dann ist das
        Kontrollkästchen gesperrt).
  \item Wirkung: Unsignierte Plugins und Plugins mit ungültiger Signatur werden geladen –
        direkt aus ihrem Ordner, nicht aus einer geprüften Kopie. Änderungen am Code wirken nach
        \menu{Plugins,Reload plugins}.
  \item Kennzeichnung: Diese Plugins heißen überall \textit{Name}~\gui{[unsigned]} bzw.\
        \textit{Name}~\gui{[signature invalid]}; \gui{Problems:} meldet
        \code{… – loaded because developer mode is on}. Rechts in der Statuszeile steht dauerhaft
        das rote Banner \gui{DEVELOPER MODE – unsigned plugins allowed}.
  \item Unsignierte Pakete lassen sich im Entwicklermodus auch installieren.
\end{itemize}

Schalten Sie den Entwicklermodus nach dem Entwickeln wieder aus.

\subsection{Signieren (für Herausgeber)}

Signiert wird mit dem Werkzeug \code{nfcad-plugin} auf dem Rechner des Herausgebers; der geheime
Schlüssel (32 Byte, Dateirechte 0600) gelangt nie ins Repository. Nach jeder Änderung am Plugin
muss neu signiert werden. Die Befehle beschreibt Anhang~\ref{sec:programme-nfcad-plugin}.

\begin{lstlisting}
nfcad-plugin keygen ~/.raytatouille-fusion/publisher.key   # einmal
nfcad-plugin sign   mein_plugin ~/.raytatouille-fusion/publisher.key
nfcad-plugin verify mein_plugin
nfcad-plugin pack   mein_plugin          # -> <id>-<version>.nfplugin
\end{lstlisting}

Wer den geheimen Schlüssel besitzt, kann Plugins veröffentlichen, die jede Installation lädt.
Geht er verloren, müssen alle Plugins mit einem neuen Schlüssel neu signiert und eine neue Version
der Anwendung mit dem neuen öffentlichen Schlüssel ausgeliefert werden.

% ---------------------------------------------------------------------------
\section{Python-Umgebung der Plugins}\label{sec:verwaltung-python}

\produkt{} bringt kein eigenes Python mit. Alle Plugins laufen mit \emph{einem} Interpreter, den
Sie unter \menu{Plugins,Settings…} wählen (Abschnitt~\ref{sec:plugin-einstellungen}); er braucht
die Pakete aller geladenen Plugins. Die Anwendung startet jedes Plugin als

\begin{lstlisting}
python -m nfcad.host --server 127.0.0.1:<port> --plugin <ordner>
\end{lstlisting}

mit dem Paket \code{nfcad} der Anwendung im Suchpfad. Das Plugin verbindet sich über TCP auf der
lokalen Schnittstelle und weist sich mit einem zufälligen Schlüssel aus, den die Anwendung ihm in
der Umgebung übergibt; andere Programme auf dem Rechner werden abgewiesen. Große Datenmengen werden
als temporäre Dateien übergeben, die die Anwendung einliest und löscht.

Wechseln Sie den Interpreter, werden alle Plugin-Prozesse beendet und mit dem neuen Interpreter
neu gestartet.

% ---------------------------------------------------------------------------
\section{Übersichtsbilder der Blöcke (für Plugin-Autoren)}\label{sec:verwaltung-uebersicht}
\index{Uebersichtsbild@Übersichtsbild!Protokoll}\index{preview\_block@\texttt{preview\_block}}

Den Reiter \gui{Overview} im Parameterdialog eines Workflow-Blocks
(Abschnitt~\ref{sec:workflow-uebersicht}) füllt ein Plugin mit der Methode
\code{preview\_block(request)} seiner Plugin-Klasse. Sie gibt eine kleine 2D-Vektorzeichnung als
JSON zurück; die Anwendung skaliert sie auf den Reiter, färbt sie passend zur hellen oder dunklen
Oberfläche und setzt die Legende der Materialien. Das Plugin braucht dafür keine
Zeichenbibliothek (kein matplotlib); \code{nfcad.drawing} hilft beim Aufbau:

\begin{lstlisting}
from nfcad.drawing import Drawing, answer, unsupported

def preview_block(self, request):
    node = request["node"]          # mit den Werten im Dialog
    if node["type"] != "my.layer":
        return unsupported()        # der nächste Zeichner oder das Schema
    d = Drawing(node.get("name", "Layer"), unit="µm", axes=("x", "y"))
    d.material("si", "si · n = 3.88")
    d.polygon([(0, 0), (1, 0), (1, 0.2), (0, 0.2)], material="si")
    d.dimension((1.1, 0), (1.1, 0.2), "0.2")
    return answer(d)
\end{lstlisting}

\paragraph{Anfrage.} \code{node}: der Block, seine \code{params} schon mit den Werten im Dialog,
die noch nicht übernommen sind; \code{type}: sein Katalogeintrag; \code{model}: das ganze Modell
des Workflows (z.\,B.\ für die Glaskataloge der Quelle); \code{scope}: die Gruppen um den Block;
\code{locale}: die Sprache der Anwendung (z.\,B.\ \code{de\_DE}).

\paragraph{Antwort.} \code{\{"ok": true, "drawing": …\}} zeigt die Zeichnung;
\code{\{"ok": false, "error": "…"\}} zeigt den Fehler rot über dem Schema des Blocks;
\code{\{"ok": false, "unsupported": true\}} gibt den Block weiter. Ausnahmen werden zu Fehlern
(mit Traceback im \gui{Plugin log}); eine Antwort mit NaN oder NumPy-Zahlen wird abgewiesen.
Die Methode läuft in einem eigenen Thread, neben Aktionen und Blockauswertungen, eine Anfrage
nach der anderen; sie sollte nach höchstens einer Sekunde fertig sein.

\paragraph{Zeichnung} (Format \code{nfcad-drawing}, Version~1):

\begin{referenz}{@{}>{\raggedright\arraybackslash}p{0.2\linewidth}X@{}}
\toprule
\textbf{Schlüssel} & \textbf{Bedeutung} \\
\midrule
\endhead
\code{title} & Überschrift. \\
\code{unit}, \code{axes} & Einheit aller Koordinaten (\code{"mm"}, \code{"µm"}) und Namen der
  waagrechten und senkrechten Achse, z.\,B.\ \code{["z", "y"]}; \code{"axes": false} zeichnet
  keine Achsen (ein Schema). \\
\code{aspect} & \code{"equal"} (Vorgabe: ein Maßstab für beide Achsen) oder \code{"fit"} (jede
  Achse füllt den Reiter). \\
\code{items} & \code{polygon} (\code{points}, mindestens drei, gefüllt; \code{material},
  \code{fill} als \code{\#rrggbb}, \code{hatched}), \code{polyline} (\code{points}, mindestens
  zwei), \code{text} (\code{at}, \code{text}, \code{align} \code{left}/\code{center}/\code{right},
  \code{valign} \code{top}/\code{middle}/\code{bottom}), \code{dimension} (\code{from},
  \code{to}, \code{text}: Maßpfeil mit Wert), \code{arrow} (\code{from}, \code{to}, \code{text}). \\
\code{role} & Bedeutung einer Linie, eines Pfeils oder Texts, bestimmt Farbe und Strich:
  \code{outline} (Vorgabe), \code{surface}, \code{axis} (strichpunktiert), \code{ray},
  \code{light}, \code{stop}, \code{detector}, \code{dimension}, \code{guide} (gepunktet),
  \code{note}, \code{error}; \code{color}, \code{width} (Pixel), \code{dashed} überschreiben sie. \\
\code{materials} & Legende: \code{\{"name", "label", "fill"\}}; ein Material, das nur ein Polygon
  nennt, kommt hinzu. Ohne \code{fill} folgt die Farbe aus dem Namen (dasselbe Material sieht in
  jedem Bild gleich aus, Luft fast durchsichtig). \\
\code{notes} & Zeilen unter dem Bild. \\
\bottomrule
\end{referenz}

Koordinaten sind endliche Zahlen, $y$ nach oben; höchstens 20\,000 Elemente und 400\,000 Punkte.
Eine fehlerhafte Zeichnung weist die Anwendung mit Grund ab (rot über dem Schema), unbekannte
Elementtypen überspringt sie – ein neueres Plugin zeichnet also auch in einer älteren Anwendung.
\code{nfcad.drawing.check(json)} nennt, was die Anwendung abweisen würde (für die Tests eines
Plugins).

\paragraph{Wer gefragt wird.} Zuerst der Anbieter des Blocktyps (das Plugin im Katalogeintrag
bzw.\ unter \code{providers} des Workflows, sonst das Plugin, dessen \code{blocks()} den Typ
enthält), dann jedes Plugin, dessen \datei{plugin.json} den Typ unter \code{previews} nennt. Die
Anwendung startet das Plugin dafür, falls nötig. Antwortet keines mit einer Zeichnung, zeigt der
Dialog das Schema des Blocks. Ein künftiger Kern (z.\,B.\ FEM) zeichnet den Schnitt seines Netzes
genauso: Polygone je Material, \code{aspect} \code{equal}, \code{unit} des Netzes.

\paragraph{Nachrichten.} Anwendung an Plugin \code{\{"type": "preview", "rid", "request"\}},
Plugin an Anwendung \code{\{"type": "previewed", "rid", "response"\}}; JSON-Zeilen wie alle
anderen Nachrichten.

% ---------------------------------------------------------------------------
\section{Update-Prüfung}\label{sec:verwaltung-updates}\index{Update-Pruefung@Update-Prüfung}

\menu{Help,Check for updates…} prüft – nur wenn Sie es wählen, nie von selbst –,

\begin{itemize}
  \item ob es eine neuere Version von \produkt{} gibt (letzte Veröffentlichung auf GitHub,
        \code{PhotonPower/Raytatouille-Fusion}) und
  \item für jedes geladene Plugin mit einem Feld \code{update\_url} in seiner
        \datei{plugin.json}, ob dort eine neuere Version angeboten wird.
\end{itemize}

Die Statuszeile zeigt währenddessen \gui{Checking for updates …}. Das Ergebnis erscheint im
Fenster \gui{Updates}, je Zeile:

\begin{itemize}
  \item \textit{Name}: \textit{neu} \code{available (installed} \textit{alt}\code{) – download}
        mit einem Link zur Download-Seite bzw.\ zum Paket,
  \item \textit{Name}~\textit{Version}~\code{is up to date},
  \item \textit{Name}~\textit{Version}\code{: check failed (…)} bei einem Netzwerkfehler.
\end{itemize}

Die Adresse in \code{update\_url} muss ein JSON-Dokument der Form
\code{\{"version": "1.2.0", "url": "<Download-Seite oder .nfplugin>"\}} liefern. Die Anfragen
senden keine persönlichen Daten; als Kennung dient nur \code{RaytatouilleFusion/}\textit{Version}.
Jede Anfrage hat 15~Sekunden Zeit. Ein heruntergeladenes Paket installieren Sie anschließend mit
\menu{Plugins,Install plugin…}; die Anwendung lädt selbst nichts herunter.

Die mitgelieferten Plugins haben kein \code{update\_url}; sie werden mit neuen Versionen der
Anwendung aktualisiert.
